\clearpage
\VerdierPhysicalPageAnchor{0178}%
\begingroup\small\setlength{\parskip}{0pt}\setlength{\emergencystretch}{1.5em}
\setlength{\abovedisplayskip}{4pt plus 2pt minus 2pt}%
\setlength{\belowdisplayskip}{4pt plus 2pt minus 2pt}%

\VerdierSemanticAnchor{ch3d-1-2-3-a-cont}
\noindent\textit{は厳密に充満忠実である。これらは対象上単射であり、したがって
圏 $K^{,*}(\mathcal A)/\operatorname{Ac}(\mathcal A)$ を
$D(\mathcal A)$ の厳密充満三角部分圏として実現する。$D(\mathcal A)$ の
対象 $X$ が $K^{,b}(\mathcal A)/\operatorname{Ac}(\mathcal A)$
（それぞれ $K^{,+}(\mathcal A)/\operatorname{Ac}(\mathcal A)$、
それぞれ $K^{,-}(\mathcal A)/\operatorname{Ac}(\mathcal A)$）の対象である
ための必要十分条件は、$|i|$ が大きいとき（それぞれ $i$ が小さいとき、
それぞれ $i$ が大きいとき）$H^i(X)=0$ となることである。}

\VerdierSemanticAnchor{ch3d-1-2-3-b-diagrams}
\noindent\textit{b) 次の図式に実線で現れる自然関手は：}
\begin{center}
\resizebox{0.80\textwidth}{!}{%
$\begin{tikzcd}[ampersand replacement=\&,column sep=4.2em,row sep=2.7em]
K^{+}(\mathcal A)/\operatorname{Ac}^{+}(\mathcal A)
  \arrow[r] \& K^{,+}(\mathcal A)/\operatorname{Ac}(\mathcal A)
  \arrow[dr,dashed] \& \\
\& \& D(\mathcal A) \\
K^{-}(\mathcal A)/\operatorname{Ac}^{-}(\mathcal A)
  \arrow[r] \& K^{,-}(\mathcal A)/\operatorname{Ac}(\mathcal A)
  \arrow[ur,dashed] \&
\end{tikzcd}$%
}

\medskip
\resizebox{0.88\textwidth}{!}{%
$\begin{tikzcd}[ampersand replacement=\&,column sep=3.2em,row sep=2.1em]
\& K^{+,b}(\mathcal A)/\operatorname{Ac}^{+}(\mathcal A) \arrow[dr] \& \& \\
K^{b}(\mathcal A)/\operatorname{Ac}^{b}(\mathcal A)
  \arrow[ur] \arrow[dr]
  \& \& K^{,b}(\mathcal A)/\operatorname{Ac}(\mathcal A)
  \arrow[r,dashed] \& D(\mathcal A) \\
\& K^{-,b}(\mathcal A)/\operatorname{Ac}^{-}(\mathcal A) \arrow[ur] \& \&
\end{tikzcd}$%
}
\end{center}
\noindent\textit{圏同値である。これらは対象の集合上で単射である。}

\VerdierSemanticAnchor{ch3d-1-2-3-c-diagram}
\noindent\textit{c) 次の図式に実線で現れる自然関手は：}
\begin{center}
\resizebox{0.82\textwidth}{!}{%
$\begin{tikzcd}[ampersand replacement=\&,column sep=4.0em,row sep=2.7em]
K^{+,b}(\mathcal A)/\operatorname{Ac}^{+}(\mathcal A)
  \arrow[r] \& K^{+}(\mathcal A)/\operatorname{Ac}^{+}(\mathcal A)
  \arrow[dr,dashed] \& \\
\& \& D(\mathcal A) \\
K^{-,b}(\mathcal A)/\operatorname{Ac}^{-}(\mathcal A)
  \arrow[r] \& K^{-}(\mathcal A)/\operatorname{Ac}^{-}(\mathcal A)
  \arrow[ur,dashed] \&
\end{tikzcd}$%
}
\end{center}
\noindent\textit{厳密に充満忠実であり、対象の集合上で単射である。}

\VerdierSemanticAnchor{ch3d-1-2-3-proof}
これは (1.2.1) および（第 II 章，2.3.5）から直ちに従う。
\endgroup
